\expandafter\ifx\csname ifdraft\endcsname\relax
   \documentclass[a4j,12pt]{jsreport}
% 数式
\usepackage{amsmath,amsfonts}
\usepackage{bm}

% 画像
\usepackage[dvipdfmx]{graphicx}

% biblatex
\usepackage[style=numeric, sorting=none, maxbibnames=1000]{biblatex}
\addbibresource{graduation_thesis.bib}

%itembox
\usepackage{ascmac}
%breakitemboxのためのpackage
\usepackage{itembkbx}
% \usepackage{tcolorbox} 画像が見えなくなってしまう原因
% \usepackage{framed} 枠

%フォント
\usepackage[ipaex]{pxchfon}

\usepackage[top=35truemm,bottom=30truemm,left=30truemm,right=30truemm]{geometry}
\usepackage{float}
\usepackage{multicol}
\usepackage{paralist}
\usepackage{caption}
%\captionsetup{format=plain, justification=centering}
\usepackage{setspace}
\usepackage{algorithm}
\usepackage{algpseudocode}
\usepackage[dvipdfmx,hidelinks]{hyperref}
\usepackage{pxjahyper}
\usepackage{fancyhdr}
\usepackage{ltablex}
\usepackage{capt-of}
\usepackage{tabularx}
\usepackage{enumitem} % itemizeの行間 これは下になければならない
\usepackage{threeparttable} % 表の脚注


   \setlength{\headheight}{16.5pt}
   \captionsetup{format=hang}
   \begin{document}
   % 上部のヘッダー
\pagestyle{fancy}
\lhead{\rightmark}
\renewcommand{\chaptermark}[1]{\markboth{第\ \thechapter\ 章\ #1}{}}
\fi

\chapter{システムに搭載したデータ}\label{appendix:story_structure}
ここでは, 本システムに搭載した『ブラック・ジャック』の物語構造分析から得られたデータを示す.
\section{設定可能なジャンル}\label{appendix:genre_list}
\begin{table}[H]
   \centering
   \caption{ジャンル一覧}
   \label{genre_list}
   \begin{tabular}{cccccc}
      \hline
      学園 & 自然 & 生き物 & 公害 & 戦争 & ドラマ \\ 
      政治 & SF & バイオレンス & 恋愛 & 真相 & 冒険 \\ 
      学生 & 放浪 & 家族 & 犯罪 & 真実 & 愛情 \\ 
      軍事 & 事件 & 医療 & 機械 & 動物 & 医学 \\ 
      友情 & 学業 & 芸術表現 & 海外 & 悲劇 & 研究 \\ 
      サスペンス & 学生生活 & 災害 & ミステリー & ヤクザ & 人間ドラマ \\ 
      道徳 & スポーツ & 挑戦 & 倫理 & トラウマ & 学校生活 \\ 
      旅 & 歴史 & 環境問題 & 事故、災害 & 妄想 & 教育 \\ 
      青春 & 過去 & ホラー & スリリング & 争い & 正義 \\ 
      学園生活 & 事故 & 善悪 & モラル & 兄弟 & 事故・災害 \\ 
      宗教 & 愛憎劇 & 復讐劇 & 植物 &  &  \\
      \hline
   \end{tabular}
\end{table}

\newpage

\section{設定可能なキャラクター属性}\label{appendix:chara_attribute_list}

\setlongtables
\begin{tabularx}{\textwidth}{lXp{1cm}p{2.5cm}}
   \caption{既存キャラクターの設定（その1） \label{chara_setting_list1}}
   \vspace*{-5pt}
   \endfirsthead
   \hline
   キャラクター名 & 役割 & 性別 & 年齢 \\
   \hline
   ブラック・ジャック & 主人公, 協力者, 敵対者, 援助者, 犠牲者, 依頼者, ライバル & 男 & 30代, 10代 \\
   ピノコ & 主人公, 協力者, 敵対者, 援助者, 犠牲者, 依頼者, ライバル & 女 & 18歳, 19歳, 21歳, 20歳 \\
   本間丈太郎 & 主人公, 協力者, 敵対者, 援助者, 犠牲者, 依頼者, ライバル & 男 & 故人, 老人 \\
   如月恵 & 主人公, 協力者, 敵対者, 援助者, 犠牲者, 依頼者, ライバル & 女 & 30代 \\
   ドクター・キリコ & 主人公, 協力者, 敵対者, 援助者, 犠牲者, 依頼者, ライバル & 男 & 30代 \\
   \hline
\end{tabularx}

% キャラクター名, 身分・職業, 特徴, 性格の表
\setlongtables
\begin{tabularx}{\textwidth}{lXX}
   \caption{既存キャラクターの設定（その2） \label{chara_setting_list2}} \vspace*{-5pt}
   \endfirsthead
   \hline
   キャラクター名 & 身分・職業 & 特徴 \\
   \hline
   ブラック・ジャック & 医者, 学生 & アウトロー, インターン, 患者, 助手\\
   ピノコ & 助手, 同居人, 患者 & 子ども, 妻, 患者 \\
   本間丈太郎 & 元外科医, 医者 & 名医, 恩師, 世界一の外科医, 故人 \\
   如月恵 & 船医 & 元恋人 \\
   ドクター・キリコ & 医者 & 安楽死, 息子 \\
   \hline
\end{tabularx}


\setlongtables
\begin{tabularx}{\textwidth}{lX}
   \caption{既存キャラクターの設定（その3） \label{chara_setting_list3}}
   \vspace*{-5pt}
   \endfirsthead
   \hline
   キャラクター名 & 性格 \\
   \hline
   ブラック・ジャック & 熱心, 早とちり, 協力的, 強い信念, 実直, 頼り甲斐がある, 思いやりがある, 我慢強い, 厳しい, 信念がある, 意思が強い, 優しい, 面倒見がよい, 負けず嫌い, 人情家, 淡白, 冷酷, 弱気, 恩師思い, 執念深い, 母親思い, 母思い, 賢い, 頑な, 平静, 落ち着いている \\
   ピノコ & わがまま, ヒステリー, ミーハー, 一途, 人騒がせ, 健気, 働き者, 優しい, 勇敢, 動物好き, 強気, 怖いもの知らず, 思いやりがある, 悪役好き, 頑張り屋, 頼りになる, 明るい, 淡白, 早とちり, 社交的, 根性がある, 面倒見がいい, 焼きもち焼き, 真面目, 涙もろい, 素直, 能天気 \\
   本間丈太郎 & 頼り甲斐がある, 信念がある, 厳しい, 現実的 \\
   如月恵 & 優しい, 前向き \\
   ドクター・キリコ & 合理的, 冷静, 頑固, 協力的, 早とちり \\
   \hline
\end{tabularx}



\setlongtables
\begin{tabularx}{\textwidth}{l|X}
   \caption{新規キャラクターの属性 \label{chara_setting_list4}}
   \vspace*{-5pt}
   \endfirsthead
   \hline
   属性の名称 & 要素  \\
   \hline
   身分・職業 & ライオン, 猿, 犬, 動物, クマ, サボテン, 大木, 患者, 病原菌, 医者, 猫, 競走馬, シカ, 鳥, 山猫, プロミンスの恋人, 鯉, さよりのペット, 夫人, 身分が高い, お金持ち, 貿易商, 処理班, 宇宙人, ー, 研究者, 故人, 大親分, 胎児, 学生, 高校退学, 中学生, タレント, 高校二年生, 社会人, 村人, 高校生, 姫, 少年船員, ピアニスト, 強盗, 殺人犯, 化石発掘, 羊毛会社社長, 軍人, 息子, 島民, ミクチャ族, 甲板員, 被害者, 看護師, 看護婦, 漫画家, ウェイトレス, 妹, 女優, 重症患者, 大地主, 元ホステス, ヤクザの女, 会社員, 売店スタッフ, 幼稚園の先生, 歌手, 配膳, 妻, アニメーター, 博士, 医学生, 跡取り, 研修医, 代役, カーレーサー, 犯罪者, 兄, 大工, ヤクザ, プロ野球選手, 盗賊, タクシー運転手, プロミンスの馬主, 仕立屋, 俳優, 画家, 金持ちの息子, 夫, 医者の妻, 後妻, 祈祷師, 母親, 内科医, 主婦, 元妻, 女医, 妊婦, 裁判長, 大統領夫人, 教師, 先生, 人さらい, 科学アカデミー副会長, 人妻, 住民, 院長, 犯罪組織, 犯罪組織所属, 犯罪組織の首領, 部下, 猟師, 警部, 警察官, 商社, 船長, 船乗り, 建設会社, 外科医, 頭取, 観光客, 偽物, 家主, サラリーマン, 暴力団, 土木, 浮気相手, 暗殺者, 運送会社, 小学校教師, 運転手, 料理人, 救急病院の医師, 教員, 悪党, 役人, クリエイター, 管理, 専務, 密猟者, 容疑者, 技師, 営業部長, 協会運営, 同級生, 召使, 土木作業員, 借金取り, 保険医, 出版社勤務, 若旦那, 暴力団員, 桑田このみの恋人, 皇帝陛下, 牧童頭, 泥棒, 土木技師, 幼稚園バスの運転手, 超能力者, イングリッドの父親, テロリスト, 寿司職人, 自然保護団体, 資産家, 伯爵, 手下, 空軍少佐, 賭博組織, 社長, 渡し守, 看護, 姐御, 助手, \\
\end{tabularx} 

\setlongtables
\begin{tabularx}{\textwidth}{l|X}
   　　　　　&芸能マネージャー, パリ市警察局警部, 誘拐犯, 印刷業, 刑事, 父親, ぺろの飼い主, 農家, 工場長, 土建屋,  花火工場, 建設業者, 五増緒部屋, 政府関係者, 大統領, 大僧官, 外務大臣, サーカス団長, 弁護士, 工事作業員, プロレス協会, 学校の先生, 下男, 権力者, 学者, マネージャー, パイロット, 僧侶, ギャング, 農園主, 石油王, 地主, 牧場主, 石油会社, 元教師, 映画監督, 印刷工場, 神父, ニッパチ産業社長, 配給会社, 町人, ハリ師, 代議士, 教授, 料理店, 家元, 校長, 医師, 獣医, 元現場責任者, スリ師, 高利貸し, 園長, 心霊研究家, 不動産, 医師連盟会長, 車掌, プロデューサー, クワン共和国大統領, 子ども, 民宿, 小説家, 霊媒師, 王, 事業家, マフィア, 土地の所有者, 甥, 郵便局長, 娘, ハイネンマンの娘, 小学生, 孤児, 児童, 県会議員の息子, 大統領の息子, 幼稚園児, 小学六年生, 学園の理事, 副社長, カイナン民主共和国大使, バイオリニスト, コロンビア大学教授, 医官, 甚大の妻, 刀鍛冶, 中華料理店経営, トキワ荘の大家, 実業家, 科学アカデミー名誉教授\\
   \hline
   特徴 & 白い, 小さい, 狂犬病, 宿主, 人殺し, オプンチア, ケヤキの木, 天然記念物, 水分, コンピューター, 恩返し, 野良猫, 馬, 肥大した脳, シャチ, 孤独, 賢い, 母親, 危険を察知する, 錦鯉, 人間の声, 死亡, 妊婦, 超能力, ピノコの姉, 畸形嚢腫, 父親, お金持ち, 極秘任務, 夫, 行方不明, ロシア人, 元患者, イタリア人, 新生児, 娘, 生命力, 虚弱体質, 不良, 姉, 患者, 放送部, 新人, 鳥人間, 働き者, 面倒見がよい, 重篤, 重病人, 先祖思い, 同級生, 晦日市病, 不良グループのリーダー, 医学部一年生, 長男, 自殺未遂, はみだしっ子, 息子, 素直, 日本好き, お笑い好き, 記憶喪失, 水泳選手, 貧乏, 難病, 怪我人, 兄思い, 弟, 恩人, 受験生, 殺人犯, 養鯉家, たくましい, 柔道部主将, SF好き, 友人, 故人, ー, 秀才, 剣道部, 推理作家志望, デベソ, 強がり, 一卵性双生児, 孤児, 小児麻痺, ニュージーランド人, 親分, 重傷者, 車椅子, スリ, 笑い声, 恋人, 野生, ノイローゼ, 延命, 交通事故, クル病, 正直, 身分が高い, 本間丈太郎の娘, 花嫁, 勇敢, 親切, 養女, 幽霊, 整形, 重体, 赤い痣, 整形手術, 殺人未遂, ガン, 若い, 姐御, 騙されやすい, ヒステリー, 平等, 割れた口, 患者思い, ミイラ, 妊娠六ヶ月, 純粋, 義手, 浮気者, アメリカ留学中, 社長の息子, 病院勤務, ブサイク, 空腹, 自分本位, 手足の痺れ, 心臓病, 死刑囚, 顔がいい, 盗賊, 弟子, 病人, 被害者, 子分, 注目選手, 正義感, 革命グループ, 仲間, レジスタンス, 手術跡, 若だんな, 親孝行, 怪人ゴーラス役, 犯罪者, 外国人, リーダー, 元教師,\\
\end{tabularx} 
   
  
\setlongtables
\begin{tabularx}{\textwidth}{l|X}
   　　　　　&  子どもがいない, 熱心, 妹, マカオ人, 原告, 教祖, 信心深い,生徒思い, 哲男の母親, 植物人間, 背骨骨折, 冷酷, 明の妻, 肺塞栓, 妻, 色盲, 自殺願望, 中毒症状, 逃亡, 息子思い, Rh-型, 技術者, ドイツ人, フィリピン人, 元恋人, 助手, 隣人, いい加減, 人任せ, エゴ, 部下, 間久部緑郎の部下, 国際指名手配, 脅迫, 有資格, 兄, 捜査, 経理課長, 操縦士, 食道ガン, 競馬好き, 社員, 世界的な外科医, 外科医長, 白人, 丸徳銀行勤務, 兄弟思い, 指名手配, 計画性がある, 頭が悪い, アラビア人, 主治医, 連続犯, 哲男の父親, 共犯者, プロフェッショナル, ひき逃げ, 言葉が悪い, 大きな手, 反対派, 旦那, 礼儀正しい, 担任, 町の医者, 回生病院勤務, ボス, D国出身, 記憶障害, 多才, 九州出身, 爆破犯人, 上役, 組員, ノーベル医学賞受賞, 四男, 嘱託医, 気が短い, ブラック・ジャックの同級生, 幹事, 外科医, 救急病院勤務, 大火傷, 弟思い, サラ金, 調査員, 幼なじみ, 洋服店経営, 運び屋, 婚約者, 医者, 家畜の世話, 犯罪グループ, 単独行動, カトリック教徒, 黒人, スウェーデン人, 安楽死, 脱獄囚, 覆面, 両腕がない, 混血児, アメリカ人, 大財閥, 特殊訓練軍, 幹部, 芸能プロダクション, 体調不良, 失踪, 内科医, 院長, 局長, 部長, 捜査中, 助教授, 社長, 手術料未払い, お節介, 経験不足, 弱気, 労災病院医師, 飲んだくれ, 狡猾, 出稼ぎ, 小説のモデル, 麻薬売買, 三下, 親方, 元旦那, 独裁者, 嘘つき, 父親代わり, 捜査三課, マタギ, 突平の弁護士, 秘密保持者, 依頼者, 副院長, 会長, 江戸っ子, 親身, 元自衛官, 外科部長, おしゃべり, 顔役, 仕事熱心, 仲介役, 爆弾魔, 重症, 元兵士, 元ゲリラ軍, アルコール中毒, 有力者, フランス人, 短指症, 遺族, 植物状態, 敵, アフリカ人, オーストラリア人, 無免許, 精神科, 名医, 名監督, まとめ役, 評判悪い, 権力者, おじ, 盲目, 医長, 心霊超科学研究者, 無欲, 逮捕者, 自然保護, 経験豊富, 恩師, X大病院勤務, 動物思い, がめつい, 責任者, 元名誉教授, 考古学研究, 町医者, 祖父, 知り合い, 金儲け, 音楽好き, アラブ人, 悪名たかい, 内臓全転位症, シャムの双生児, 村人, 貪欲, 実力がある, 世話好き, マカオ在住, 金目当て, 肥満, 大スター, 住人, 白血病, 山研究家, 関根の妻, 沖縄出身, 校医, 父親思い, ハセドー病, 針が怖い, ピノコのモデル, 甥, ベトナム人, 医者の息子, ヤクザの子ども, 強引, 奇形児, 川崎病, 仮病, 孫, いたずらっ子, エディプスコンプレックス, 判断力がある, 役者, 魚鱗癬, そろばんが得意, 博の恋人, 火 悪性貧血, 無頭児, インテリ, 競輪好き, 大金持ち, 重役, アフリカ在住, 大先生, 心臓外科の第一人者, 秘密厳守,  六つ子,\\
\end{tabularx} 

    
\setlongtables
\begin{tabularx}{\textwidth}{l|X}
   　　　　　& 守秘義務, ストラディバリウス, トミーの父親, アウトロー, 有名外科医, 高位, 青い瞳, とび色の瞳, 衰弱, 未亡人, 技術力がある, 中国人\\
   \hline
   性格 & 忠実, 凶暴, 荒い, 寂しがり, 優しい, 野生的, 賢い, 誠実, 甘えん坊, 勇敢, 子ども思い, 健気, ぐうたら, 利口, 放任主義, 気が強い, 自分本位, 前向き, \\
   &妻思い, 薄情者, あくどい, 心優しい, 大人しい, ネガティブ, 臆病, 根気強い, 強気, 母親思い, 友だち思い, 情熱的, おしとやか, 勇気がある, 人懐っこい, 子ども好き, 冷酷, 悲観的, 弱気, 動物好き, 好奇心旺盛, 朗らか, ポジティブ, 無鉄砲, 努力家, 負けず嫌い, 意地っ張り, 父親思い, 心配性, 思いやりがある, 真面目, 素直, 生意気, 盲目的, 自分勝手, 正義感が強い, 嫉妬深い, 積極的, 暴力的, 熱心, 頑固, 純粋, 愛情深い, 戦争嫌い, 明るい, 思い込みが激しい, 頑張り屋, 責任感がない, 荒っぽい, わがまま, 懸命, 我慢強い, 物怖じしない, 兄思い, 温和, 人がいい, 大袈裟, 責任感が強い, 単純, 弟思い, 親思い, 頼り甲斐がある, 自信家, 気が短い, 恋人思い, 猟奇的, いい加減, 根性がない, 必死, 諦めが悪い, 強欲, 執念深い, 無謀, がめつい, 押し付けがましい, 息子思い, 厳しい, 献身的, はっきりしている, 包容力がある, しっかりしている, 無責任, 研究熱心, したたか, 怠け者, 残忍, 嘘つき, 強引, 熱血, 冷静, 実直, 怒りっぽい, 失礼, プライドが高い, 自己中心的, 慎重, 人騒がせ, 情深い, よそよそしい, 調子が良い, ずる賢い, おっちょこちょい, 義理堅い, 頼りになる, 勝気, 楽観的, 協力的, 筋が通っている, 卑怯, 判断力がある, 現実的, 冷徹, 後先考えない, ケチ, しつこい, 非協力的, 協調性がない, 理解力がある, 合理的, 一途, 残酷, 勘がいい, 自暴自棄, 衝動的, 仕事熱心, のんびり屋, ワンマン, 不誠実, 決断力がある, 図々しい, 傲慢, 見栄っ張り, 娘思い, せっかち, 押しが強い, 機械的, 落ち着きがない, 疑い深い, 大雑把, 正直, 正義感がある, 好戦的, 信念がある, 利己的, 行動力がある, 家族思い, 動物思い, 丁寧, 愚痴っぽい, 一生懸命, 面倒見がいい, 礼儀正しい, 横柄, 騒がしい, 芯がある, 孫思い, 変人, 繊細, 世話焼き, 怖がり, 意地悪, 妹思い, 鳥思い, 国民思い, 人情家, 人任せ, 用心深い, 貪欲, 責任感がある, 意思が強い, 生命力がある, 変わり者, 身勝手,\\
   \hline
\end{tabularx} 
% \begin{breakitembox}
%    [l]{キャラクター属性における身分・職業の一覧}
%    ライオン, 猿, 犬, 動物, クマ, サボテン, 大木, 患者, 病原菌, 医者, 猫, 競走馬, シカ, 鳥, 山猫, プロミンスの恋人, 鯉, さよりのペット, 夫人, 身分が高い, お金持ち, 貿易商, 処理班, 宇宙人, ー, 研究者, 故人, 大親分, 胎児, 学生, 高校退学, 中学生, タレント, 高校二年生, 社会人, 村人, 高校生, 姫, 少年船員, ピアニスト, 強盗, 殺人犯, 化石発掘, 羊毛会社社長, 軍人, 息子, 島民, ミクチャ族, 甲板員, 被害者, 看護師, 看護婦, 漫画家, ウェイトレス, 妹, 女優, 重症患者, 大地主, 元ホステス, ヤクザの女, 会社員, 売店スタッフ, 幼稚園の先生, 歌手, 配膳, 妻, アニメーター, 博士, 医学生, 跡取り, 研修医, 代役, カーレーサー, 犯罪者, 兄, 大工, ヤクザ, プロ野球選手, 盗賊, タクシー運転手, プロミンスの馬主, 仕立屋, 俳優, 画家, 金持ちの息子, 夫, 医者の妻, 後妻, 祈祷師, 母親, 内科医, 主婦, 元妻, 女医, 妊婦, 裁判長, 大統領夫人, 教師, 先生, 人さらい, 科学アカデミー副会長, 人妻, 住民, 院長, 犯罪組織, 犯罪組織所属, 犯罪組織の首領, 部下, 猟師, 警部, 警察官, 商社, 船長, 船乗り, 建設会社, 外科医, 頭取, 観光客, 偽物, 家主, サラリーマン, 暴力団, 土木, 浮気相手, 暗殺者, 運送会社, 小学校教師, 運転手, 料理人, 救急病院の医師, 教員, 悪党, 役人, クリエイター, 管理, 専務, 密猟者, 容疑者, 技師, 営業部長, 協会運営, 同級生, 召使, 土木作業員, 借金取り, 保険医, 出版社勤務, 若旦那, 暴力団員, 桑田このみの恋人, 皇帝陛下, 牧童頭, 泥棒, 土木技師, 幼稚園バスの運転手, 超能力者, イングリッドの父親, テロリスト, 寿司職人, 自然保護団体, 資産家, 伯爵, 手下, 空軍少佐, 賭博組織, 社長, 渡し守, 看護, 姐御, 助手, 芸能マネージャー, パリ市警察局警部, 誘拐犯, 印刷業, 刑事, 父親, ぺろの飼い主, 農家, 工場長, 土建屋, 花火工場, 建設業者, 五増緒部屋, 政府関係者, 大統領, 大僧官, 外務大臣, サーカス団長, 弁護士, 工事作業員, プロレス協会, 学校の先生, 下男, 権力者, 学者, マネージャー, パイロット, 僧侶, ギャング, 農園主, 石油王, 地主, 牧場主, 石油会社, 元教師, 映画監督, 印刷工場, 神父, ニッパチ産業社長, 配給会社, 町人, ハリ師, 代議士, 教授, 料理店, 家元, 校長, 医師, 獣医, 元現場責任者, スリ師, 高利貸し, 園長, 心霊研究家, 不動産, 医師連盟会長, 車掌, プロデューサー, クワン共和国大統領, 子ども, 民宿, 小説家, 霊媒師, 王, 事業家, マフィア, 土地の所有者, 甥, 郵便局長, 娘, ハイネンマンの娘, 小学生, 孤児, 児童, 県会議員の息子, 大統領の息子, 幼稚園児, 小学六年生, 学園の理事, 副社長, カイナン民主共和国大使, バイオリニスト, コロンビア大学教授, 医官, 甚大の妻, 刀鍛冶, 中華料理店経営, トキワ荘の大家, 実業家, 科学アカデミー名誉教授
%    \label{chara_ocupation_list}
% \end{breakitembox}

% \vspace{1cm}

% \begin{breakitembox}
%    [l]{キャラクター属性における特徴の一覧}
%    白い, 小さい, 狂犬病, 宿主, 人殺し, オプンチア, ケヤキの木, 天然記念物, 水分, コンピューター, 恩返し, 野良猫, 馬, 肥大した脳, シャチ, 孤独, 賢い, 母親, 危険を察知する, 錦鯉, 人間の声, 死亡, 妊婦, 超能力, ピノコの姉, 畸形嚢腫, 父親, お金持ち, 極秘任務, 夫, 行方不明, ロシア人, 元患者, イタリア人, 新生児, 娘, 生命力, 虚弱体質, 不良, 姉, 患者, 放送部, 新人, 鳥人間, 働き者, 面倒見がよい, 重篤, 重病人, 先祖思い, 同級生, 晦日市病, 不良グループのリーダー, 医学部一年生, 長男, 自殺未遂, はみだしっ子, 息子, 素直, 日本好き, お笑い好き, 記憶喪失, 水泳選手, 貧乏, 難病, 怪我人, 兄思い, 弟, 恩人, 受験生, 殺人犯, 養鯉家, たくましい, 柔道部主将, SF好き, 友人, 故人, ー, 秀才, 剣道部, 推理作家志望, デベソ, 強がり, 一卵性双生児, 孤児, 小児麻痺, ニュージーランド人, 親分, 重傷者, 車椅子, スリ, 笑い声, 恋人, 野生, ノイローゼ, 延命, 交通事故, クル病, 正直, 身分が高い, 本間丈太郎の娘, 花嫁, 勇敢, 親切, 養女, 幽霊, 整形, 重体, 赤い痣, 整形手術, 殺人未遂, ガン, 若い, 姐御, 騙されやすい, ヒステリー, 平等, 割れた口, 患者思い, ミイラ, 妊娠六ヶ月, 純粋, 義手, 浮気者, アメリカ留学中, 社長の息子, 病院勤務, ブサイク, 空腹, 自分本位, 手足の痺れ, 心臓病, 死刑囚, 顔がいい, 盗賊, 弟子, 病人, 被害者, 子分, 注目選手, 正義感, 革命グループ, 仲間, レジスタンス, 手術跡, 若だんな, 親孝行, 怪人ゴーラス役, 犯罪者, 外国人, リーダー, 元教師, 子どもがいない, 熱心, 妹, マカオ人, 原告, 教祖, 信心深い, 生徒思い, 哲男の母親, 植物人間, 背骨骨折, 冷酷, 明の妻, 肺塞栓, 妻, 色盲, 自殺願望, 中毒症状, 逃亡, 息子思い, Rh-型, 技術者, ドイツ人, フィリピン人, 元恋人, 助手, 隣人, いい加減, 人任せ, エゴ, 部下, 間久部緑郎の部下, 国際指名手配, 脅迫, 有資格, 兄, 捜査, 経理課長, 操縦士, 食道ガン, 競馬好き, 社員, 世界的な外科医, 外科医長, 白人, 丸徳銀行勤務, 兄弟思い, 指名手配, 計画性がある, 頭が悪い, アラビア人, 主治医, 連続犯, 哲男の父親, 共犯者, プロフェッショナル, ひき逃げ, 言葉が悪い, 大きな手, 反対派, 旦那, 礼儀正しい, 担任, 町の医者, 回生病院勤務, ボス, D国出身, 記憶障害, 多才, 九州出身, 爆破犯人, 上役, 組員, ノーベル医学賞受賞, 四男, 嘱託医, 気が短い, ブラック・ジャックの同級生, 幹事, 外科医, 救急病院勤務, 大火傷, 弟思い, サラ金, 調査員, 幼なじみ, 洋服店経営, 運び屋, 婚約者, 医者, 家畜の世話, 犯罪グループ, 単独行動, カトリック教徒, 黒人, スウェーデン人, 安楽死, 脱獄囚, 覆面, 両腕がない, 混血児, アメリカ人, 大財閥, 特殊訓練軍, 幹部, 芸能プロダクション, 体調不良, 失踪, 内科医, 院長, 局長, 部長, 捜査中, 助教授, 社長, 手術料未払い, お節介, 経験不足, 弱気, 労災病院医師, 飲んだくれ, 狡猾, 出稼ぎ, 小説のモデル, 麻薬売買, 三下, 親方, 元旦那, 独裁者, 嘘つき, 父親代わり, 捜査三課, マタギ, 突平の弁護士, 秘密保持者, 依頼者, 副院長, 会長, 江戸っ子, 親身, 元自衛官, 外科部長, おしゃべり, 顔役, 仕事熱心, 仲介役, 爆弾魔, 重症, 元兵士, 元ゲリラ軍, アルコール中毒, 有力者, フランス人, 短指症, 遺族, 植物状態, 敵, アフリカ人, オーストラリア人, 無免許, 精神科, 名医, 名監督, まとめ役, 評判悪い, 権力者, おじ, 盲目, 医長, 心霊超科学研究者, 無欲, 逮捕者, 自然保護, 経験豊富, 恩師, X大病院勤務, 動物思い, がめつい, 責任者, 元名誉教授, 考古学研究, 町医者, 祖父, 知り合い, 金儲け, 音楽好き, アラブ人, 悪名たかい, 内臓全転位症, シャムの双生児, 村人, 貪欲, 実力がある, 世話好き, マカオ在住, 金目当て, 肥満, 大スター, 住人, 白血病, 火山研究家, 関根の妻, 沖縄出身, 校医, 父親思い, ハセドー病, 針が怖い, ピノコのモデル, 甥, ベトナム人, 医者の息子, ヤクザの子ども, 強引, 奇形児, 川崎病, 仮病, 孫, いたずらっ子, エディプスコンプレックス, 判断力がある, 役者, 魚鱗癬, そろばんが得意, 博の恋人, 悪性貧血, 無頭児, インテリ, 競輪好き, 大金持ち, 重役, アフリカ在住, 大先生, 心臓外科の第一人者, 秘密厳守, 守秘義務, ストラディバリウス, トミーの父親, アウトロー, 有名外科医, 高位, 青い瞳, とび色の瞳, 六つ子, 衰弱, 未亡人, 技術力がある, 中国人
%    \label{chara_feature_list}
% \end{breakitembox}

% \vspace{1cm}

% \begin{breakitembox}
%    [l]{キャラクター属性における性格の一覧}
%    忠実, 凶暴, 荒い, 寂しがり, 優しい, 野生的, 賢い, 誠実, 甘えん坊, 勇敢, 子ども思い, 健気, ぐうたら, 利口, 放任主義, 気が強い, 自分本位, 前向き, 妻思い, 薄情者, あくどい, 心優しい, 大人しい, ネガティブ, 臆病, 根気強い, 強気, 母親思い, 友だち思い, 情熱的, おしとやか, 勇気がある, 人懐っこい, 子ども好き, 冷酷, 悲観的, 弱気, 動物好き, 好奇心旺盛, 朗らか, ポジティブ, 無鉄砲, 努力家, 負けず嫌い, 意地っ張り, 父親思い, 心配性, 思いやりがある, 真面目, 素直, 生意気, 盲目的, 自分勝手, 正義感が強い, 嫉妬深い, 積極的, 暴力的, 熱心, 頑固, 純粋, 愛情深い, 戦争嫌い, 明るい, 思い込みが激しい, 頑張り屋, 責任感がない, 荒っぽい, わがまま, 懸命, 我慢強い, 物怖じしない, 兄思い, 温和, 人がいい, 大袈裟, 責任感が強い, 単純, 弟思い, 親思い, 頼り甲斐がある, 自信家, 気が短い, 恋人思い, 猟奇的, いい加減, 根性がない, 必死, 諦めが悪い, 強欲, 執念深い, 無謀, がめつい, 押し付けがましい, 息子思い, 厳しい, 献身的, はっきりしている, 包容力がある, しっかりしている, 無責任, 研究熱心, したたか, 怠け者, 残忍, 嘘つき, 強引, 熱血, 冷静, 実直, 怒りっぽい, 失礼, プライドが高い, 自己中心的, 慎重, 人騒がせ, 情深い, よそよそしい, 調子が良い, ずる賢い, おっちょこちょい, 義理堅い, 頼りになる, 勝気, 楽観的, 協力的, 筋が通っている, 卑怯, 判断力がある, 現実的, 冷徹, 後先考えない, ケチ, しつこい, 非協力的, 協調性がない, 理解力がある, 合理的, 一途, 残酷, 勘がいい, 自暴自棄, 衝動的, 仕事熱心, のんびり屋, ワンマン, 不誠実, 決断力がある, 図々しい, 傲慢, 見栄っ張り, 娘思い, せっかち, 押しが強い, 機械的, 落ち着きがない, 疑い深い, 大雑把, 正直, 正義感がある, 好戦的, 信念がある, 利己的, 行動力がある, 家族思い, 動物思い, 丁寧, 愚痴っぽい, 一生懸命, 面倒見がいい, 礼儀正しい, 横柄, 騒がしい, 芯がある, 孫思い, 変人, 繊細, 世話焼き, 怖がり, 意地悪, 妹思い, 鳥思い, 国民思い, 人情家, 人任せ, 用心深い, 貪欲, 責任感がある, 意思が強い, 生命力がある, 変わり者, 身勝手
%    \label{chara_personality_list}
% \end{breakitembox}

\chapter{提出されたプロット}\label{app:plot}
ここでは, 第2回制作会議に提出されたプロットを掲載する.

\vspace{10pt}
\begin{breakitembox}
   [l]{グループAが提出したプロット1}
   \noindent \{シーン1: 冒頭・手術\}\\
1. ブラック・ジャックが手術室で集中しながら手術を行っている。\\
2. 手術が終わり、ブラック・ジャックは一息つきながら達成感の持った目で、患者を見つめる。\\
3. しかし、その静けさは長く続かず、ピノコの叫びがブラックジャックを引き戻す。\\
4. ピノコが見ていたテレビでは、街中で突如ある青年が無差別大量殺人を起こしたことを伝える。その青年の名前は、道夫。彼は佐々木夫妻の子供だと後に明らかになる。\\
5. 事件は大きな衝撃を与え、謎の青年の行動が大きな話題を呼ぶ。 ブラック・ジャックは事件を見つめ、事件の背後にあるものを感じ取る\\
\\
 \{シーン2: 回想 - 依頼\}\\
1. 十数年前、研究者である佐々木夫妻の3歳になる長女が、火事に巻き込まれ怪我負わせてしまった。\\
2. 痛ましい事故を経験した佐々木夫は、これから生まれる長男（長女の弟）を遺伝子操作によって、飛び抜けた身体能力と優秀な頭脳を持つことができるように決意する。\\
3. 佐々木夫妻は、子供が未来に直面する未知の危険から生き抜くための準備として、この行為が必要なのではないかと考える。\\
4. 初めは不安げだった佐々木妻も、夫の情熱と子供の幸せを優先し、この道を選ぶことを理解し、それを支持する。\\
5. そして、二人は産まれゆく息子に遺伝子操作を施す。それが、道夫の出生の秘密である。\\
6.佐々木夫妻の元に、ある権力者（のちの総理大臣）がやってくる。その権力者は、自分のこれから産まれゆく子供への遺伝子操作も依頼する。\\
\\
 \{シーン3: 非対象者の異変の始まりと対象者の運命\}\\
1. それから十数年経ち、佐々木夫妻が遺伝子操作を施した子供たちの一部が、思春期に入ると異常な行動を取り始める。\\
2. 遺伝子操作から与えられた非凡な知力と身体能力は彼らに孤独感と圧倒的なプレッシャーをもたらしていた。彼らは周囲の人間を「バカ」と見下し、社会的弱者に対する攻撃的な態度をとるようになる。\\
3. 道夫が起こした無差別大量殺人も、それが原因だったことが判明する。\\
4. 道夫は事件後、裁判を経て死刑が宣告される。そのニュースは佐々木夫妻にとって大きなショックを与える。自分たちの子供が社会から排斥される事実と、道夫が誤って人々を傷つけたという事実が彼らに深い絶望感を与える。\\
5. しかし、それは始まりに過ぎなかった。遺伝子操作を行なった他の子供たちもまた、異常な行動を示し始める。一部は自殺し、一部は犯罪を犯し、一部は社会から孤立する。\\
6.佐々木夫妻は、道夫や他の子供たちを手術して「普通に戻してほしい」とブラックジャックに依頼する。\\
\\
 \{シーン4: 真相と迷い\}\\
1. ブラック・ジャックは懇願する佐々木夫妻と他の親たちの言葉に対し、彼らが子供たちを自分たちのエゴに巻き込んだことに対する報いだと怒る。\\
2. 「子供たちへの遺伝子操作は私が信じる医の道から外れる」と主張するブラック・ジャック。しかし、涙ながらに「ただ子供たちに幸せになってほしいんです」と懇願する親たちの言葉に、彼の心は揺さぶられる。\\
3. そんな中、道夫の死刑が延期され、更には釈放されるというニュースが流れる。それは時の総理大臣もまた、子供を佐々木夫に遺伝子操作をお願いしていた過去を持っていたため便宜が図られたのだ。\\
4. 時の総理大臣の息子もまた道夫と同じ、人々をバカにし、暴力に訴えるようになっており、総理大臣はブラック・ジャックにその暴力性を抑える手術を依頼する。\\
5. 多くの親たちからの手術の懇願。「幸せになってほしいんだな？」とブラックジャックは親たちに言い残し、手術へと向かう。\\
\\
 \{シーン5: 決断と手術\}\\
1. ブラック・ジャックは子どもたちの手術を決意する。遺伝子操作ではなく、脳を手術することで暴力を感じる箇所を切り取る、という手術だ。\\
2. ブラック・ジャックが子どもたちの手術に踏み切ったのは、彼が彼らの身に起きた不幸とその原因を深く理解していたからだ。彼自身も過去、人間の自由意志と医学の力が交錯した結果、壮絶な運命を辿ったことがあるからだ。\\
3. 手術は次々と成功し、子供たちは暴力性を失い、不自然なほどに無垢な子どもたちへと変わる。一切の暴力性を失い、天使のような微笑みを浮かべ続ける。\\
\\
 \{シーン6: 結果と再認識\}\\
1. 佐々木夫妻やその他の親たちは子供たちが無垢な存在に戻り、平和に笑い、愛に溢れた生活を送る姿を見て、涙を流し、感謝の言葉を口にする。何よりも、子供たちが犯罪を起こさないことに安堵をする。\\
2. その言葉に対し、ブラック・ジャックは苦い笑みを浮かべる。親たちが子供たちに求めた「幸せ」は何だったのか、その答えがまだ見当たらないことを、彼は察している。\\
3.  事件が収束に向かうにつれて、ブラック・ジャックは再び自身の考えを整理する。「子供たちは幸せそうになった」しかし、それが子供たち自身の意志で選んだ結果ではなかったという事実が彼の心を重くする。\\
4. エピソードが終わりに近づき、最終シーンではブラック・ジャックとピノコ、そして手術を受けた子供たちの笑顔が映し出される。不自然なほどに無垢な子どもたちと遊ぶブラックジャック。しかし、ブラック・ジャックの瞳には前途への決意が刻まれている。\\
5. ブラック・ジャックは、医学の発展と危険性、そして人間が抱く欲望に立ち向かい、次なる挑戦へと旅を続けるのであった。

\end{breakitembox}

\vspace{20pt}
\begin{breakitembox}
   [l]{グループAが提出したプロット2}
   \noindent\{シーン1: 依頼の拒否\}\\
   1. 物語の中心人物は、SNSインフルエンサーのかずやと彼の父親である国宝級の陶芸家。かずやには恋人の博がいる。\\
   2. かずやは見栄え重視のSNS生活を送りつつも、父親は「見栄えより本質」にこだわっていることが、父と息子の間に亀裂を生んでいる。\\
   3. ある夜、眠れぬままのかずやが闇投稿をしたところ、「らしくない」とフォロワーから反感を買い、バッシングを受ける。\\
   4. そんな中、博がかずやを支え続ける。\\
   \\
   \{シーン2: 依頼\}\\
   1. 翌日、かずやの体調が悪化し、その連絡が学校から父親に届けられる。不眠による発汗や倦怠感が、元気だった頃のかずやとは別人のようであった。\\
   2. 父親は息子の体調悪化に驚愕、世界中の名医に相談を持ち掛ける。しかし、その原因は見つからない。\\
   3. 病状を特定できぬ医師たちの中で、父親は医者として名高いブラックジャックに依頼する決断をする。\\
   4. 親の命懸けの状況を前にして、ブラックジャックは依頼を受け、かずやの元へ急行する。\\
   \\
   \{シーン3: 看病・治療\}\\
   1. ブラックジャックとピノコがかずやの元へ到着し、診察を開始。病名は「致死性家族性不眠症」（FFI）。眠る機能を失い、起き続けることで心身が衰弱し死に至る不治の病である。\\
   2. 万全を尽くすブラックジャックだが、父親と博は、病の治療法が世にないと心得ている。\\
   3. 父親は心細さと自分の心を落ち着かせるためにも陶芸に没入し過ぎてしまうが、その不器用な行動はかずやとの関係をさらに悪化する。\\
   4. ブラックジャックは一筋の希望として、唯一の治療法である脳の移植を父親に提案する。\\
   \\
   \{シーン4: 集団心理の罠とブラックジャックの怒り\}\\
   1. かずやの病状がSNSに広まるにつれ、フォロワーたちは深刻さに気付き始める。\\
   2. ブラックジャックが提唱する移植手術は、脳の一部を移植し睡眠機能を取り戻すものだが、犠牲者が必要となる。\\
   3. 博はかずやのSNSを勝手に更新し、フォロワーやファンに対し、かずやのための脳の提供を呼びかける。\\
   4. 集団心理と自己承認欲求に煽られ、トップフォロワーの文雄が自らの脳を提供すると宣言。その動機に、ブラックジャックは激怒。「自己承認欲求なんかで命を投げ出すな」と怒りを露わにする。\\
   5. しかし、文雄は「今のままでは僕は誰のためにもならない人生だ」と弱々しく微笑む。その微笑みに言葉を返せなくなったブラックジャックは、矛盾した感情を抱きつつも手術の準備へと進むことを決める。\\
   \\
   \{シーン5: 変化と激動の中で\}\\
   1. 手術の直前、かずやと博は父親に知らせずに手術をすることを決断をしたことをブラックジャックに打ち明けます。「本質を大事にする」父が、他者の脳を移植するといった人智を超えた行動に反対することを恐れたからです。\\
   2. ブラックジャックは複雑な心情を抱きながらも、彼らが18歳で成人であることからも、父親に知らせないことに同意します。かずやは博に自身の命と向き合える勇気をもらったことを告げ、博もまたかずやの決断を祝福します。\\
   3. 父親は息子が手術を行うことを知らずに、皮肉にも陶芸に没頭しています。しかし、その真実を知ったときの落胆と憤りは想像を絶するものでしょう。\\
   4. 手術が敢行されたことを知った父親。２人が自分を見捨ててまで手術を行ったという事実への衝撃と、治療法への違和感による怒りで揺れ動きます。しかし、博の「かずやの命はかずやがどう決断してもいいはずだ」の言葉が突き刺さります。\\
   5. 手術は成功。怒りに満ちていた父親は一転、ブラックジャックに涙交じりに感謝を伝えます。息子の命が繋がったことへの喜びは全てを凌駕し、息子が自分自身を見つめ人生の選択をした勇気を認めます。父親の心情の葛藤が描かれ、家族の新たな絆が生まれる希望を見せつつ、これから先の困難を予期させます。\\
   \\
   \{シーン6: 新たな人生の始まりと苦悩\}\\
   1. 名医ブラックジャックが提唱した移植手術の結果、かずやの体に文雄の脳が移植されると、かずやの人格は徐々に変化し始めます。文雄の記憶、感情、性格がかずやの中に生まれ、新たな存在、新たな「かずや」へと生まれ変わります。\\
   2. その新たな生命、新たな人格を持つかずやは、パソコンの画面越しに自分（文雄）が事故で死んだというニュースを目の当たりにします。\\
   3. 一方、父親は息子の変わり様に混乱し、文雄という他者の存在が息子の中に深く根を下ろしている現実をどう受け止めたらいいのか、また自分の父親としての感情や役割をどうすべきかを迷います。\\
   4. 性格が変わったかずやの投稿に対するフォロワーたちの反応はさまざまで、支持者もいれば反発する者もいて、ソーシャルメディアが揺れ動きます。その中で、文雄と合体したかずやは、自分自身と向き合い、人々の心を動かすために一体何が必要なのかを問い続けます。\\
   5. 父親はそんなかずやを前に、遂には「2人の父親になろう」という新たな決意を固めます。父親として、また一人の人間として、息子と息子のために命を投げ打った少年の人生を見守り、自分の心の中にある混乱と対峙します。\\
   6. ブラックジャックはかずやの成長と変化を静かに見守り、改めて「生命の価値とは何か、人間のアイデンティティとは何か」を問います。手塚治虫らしい、人間ドラマと深深に思索する哲学を組み合わせたエンディングになります。
\end{breakitembox}

\vspace{20pt}
\begin{breakitembox}
   [l]{グループBが提出したプロット1}
   \noindent\{シーン1: 依頼\}
   1. 革新的なアンドロイド、“HeartBeat Mark II”が突如機能停止し、発明者である川村たかしがその解決の糸口を見つけられず困惑している。彼は自力では解決できず、ブラック・ジャックに助けを求める。(川村の視点)\\
   2. ブラック・ジャックはアンドロイドを人間のような存在として観察し、その心臓部分がただの機械部品で無く、生命体を理解する鍵であることを説明する。\\
   3. ブラック・ジャックは"HeartBeat Mark II"の心臓部分を調査し、心臓を制御する電子部品に異常があることを見つけ、それが機能停止の原因であることを明らかにする。そして、足元に広がりつつある全体のエネルギー供給の不安定さを川村に伝え、アンドロイドの生命の危険性を警告する。\\
   4. ブラック・ジャックの意見に対して、ドクター・キリコが口を開き、挑発する。彼の主張は"HeartBeat Mark II"が不安定になっているならば、何らかの危険を引き起こす可能性があるというものだった。この言葉が二人の間に対立の芽を生じさせる。\\
   5. 川村はドクター・キリコの主張に言葉を失い、混乱する。しかし"HeartBeat Mark II"の生命を守るために、ブラック・ジャックに全てを委ねる決断を下す。\\
   \\
   \{シーン2: 危害\}\\
   1. "HeartBeat Mark II"の生命を延長させたいブラック・ジャックは、ピノコがドクター・キリコに捕らえられたという警告の電話を受け、葛藤する。(ブラック・ジャックの視点)\\
   2. ブラック・ジャックはドクター・キリコの冷静さに怒り、アンドロイドの再起動に全力を尽くすことを誓う。\\
   3. しかし"HeartBeat Mark II"の心臓部分の異常を修正するには手術しかないとブラック・ジャックは結論付ける。\\
   4. 川村は、そのアンドロイドが人間の生命観に対する新たな視野を提供すること、そして科学と人間性の新たな境界を模索する中でその重要性を理解する。彼はブラック・ジャックが"HeartBeat Mark II"の生命を救うために手術することを支持し、その費用をすぐに提供する。\\
   5. しかし、その手術は革新的な機械科学の知識と技術を必要とし、特に心臓部分の修復は緻密で難解な作業だと分かる。不安に感じつつも、ブラック・ジャックはアンドロイドの問題を解決しようと意志を固める。\\
   6. その間も、ピノコは自らの命を救うために、ドクター・キリコによる拘束から逃れようと全力を尽くしている。\\
   \\
   \{シーン3: 看病・治療\}\\
   1. ブラック・ジャックは川村の信頼と"HeartBeat Mark II"の心臓部分の”生命”を信じて、冷静に複雑な手術に挑む。(ブラック・ジャックの視点)\\
   2. 手術は難関だった。微細な機械構造と専門的な技術が求められる"HeartBeat Mark II"の心臓部分の手術を進めるブラック・ジャック。彼の専門知識の欠如が手術の難易度をさらに上昇させる。\\
   3. しかし、ブラック・ジャックの不屈の努力はやがて結果を伴い、心臓部分の異常は修正され、手術は成功する。\\
   4. 手術後、ブラック・ジャックがピノコを救うために、ドクター・キリコの前に現れる。直接的な対立が始まる中、ブラック・ジャックはチャンスを見つけてピノコの拘束を解き、無事に救出する。\\
   5. ブラック・ジャックの前に出現したドクター・キリコは、負けたにも関わらず、「ブラック・ジャック、今回もお前の負けだ」と言い残し、その場から立ち去る。\\
   \\
   \{シーン4: 変革\}\\
   1. ブラック・ジャックが去った後、川村は"HeartBeat Mark II"の心臓部分が生命体としての活動を続けていることを確認し、その重要性を再認識する。(川村の視点)\\
   2. その後、彼は"HeartBeat Mark II"の心臓部分を発展させ、新たなアンドロイド開発の会社を設立することを決意する。\\
   3. 川村の会社は短期間で大成功を収め、次世代のアンドロイド開発への資金を得る。\\
   4. 一方、ブラック・ジャックはピノコとともに日々を過ごしながら、ドクター・キリコの言葉を考える時間が増えていた。\\
   \\
   \{シーン5: エピローグ\}\\
   1. しかし、その川村の成功も一夜にして崩れる。"HeartBeat Mark II"プロジェクトが頓挫した報告がブラック・ジャックに届く。(ブラック・ジャックの視点)\\
   2. ブラック・ジャックは再びオペレーションルームに向かい、"HeartBeat Mark II"の手術を行う。\\
   3. アンドロイドの心臓部分が変更されていることにブラック・ジャックは驚く。その変更が心臓の働きを阻害し、手術ができない状況を作り出していた。\\
   4. その結果、川村のアンドロイド事業の失敗により、"HeartBeat Mark II"プロジェクトが崩壊する。\\
   5. ブラック・ジャックは手術の失敗を深く後悔する。そこに、ドクター・キリコの高笑いが聞こえる。
\end{breakitembox}

\vspace{20pt}
\begin{breakitembox}
   [l]{グループBが提出したプロット2}
   \noindent\{シーン1: 依頼\}\\
   1. 革新的なアンドロイド、“ HeartBeat Mark II”が突如機能停止し、発明者である川村たかしがその解決の糸口を見つけられず困惑している。彼は自力では解決できず、ブラック・ジャックに助けを求める。(川村の視点)\\
   2. ブラック・ジャックはアンドロイドを人間のような存在として観察し、その心臓部分がただの機械部品で無く、生命体を理解する鍵であることを説明する。\\
   3. ブラック・ジャックは"HeartBeat Mark II"の心臓部分を調査し、心臓を制御する電子部品に異常があることを見つけ、それが機能停止の原因であることを明らかにする。そして、足元に広がりつつある全体のエネルギー供給の不安定さを川村に伝え、アンドロイドの生命の危険性を警告する。\\
   4. ブラック・ジャックの意見に対して、ドクター・キリコが口を開き、挑発する。彼の主張は"HeartBeat Mark II"が不安定になっているならば、何らかの危険を引き起こす可能性があるというものだった。この言葉が二人の間に対立の芽を生じさせる。\\
   5. 川村はドクター・キリコの主張に言葉を失い、混乱する。しかし"HeartBeat Mark II"の生命を守るために、ブラック・ジャックに全てを委ねる決断を下す。\\
   \\
   \{シーン2: 危害\}\\
   1. "HeartBeat Mark II"の生命を延長させたいブラック・ジャックは、ピノコがドクター・キリコに捕らえられたという警告の電話を受け、葛藤する。(ブラック・ジャックの視点)\\
   2. ブラック・ジャックはドクター・キリコの冷静さに怒り、アンドロイドの再起動に全力を尽くすことを誓う。\\
   3. しかし"HeartBeat Mark II"の心臓部分の異常を修正するには手術しかないとブラック・ジャックは結論付ける。\\
   4. 川村は、そのアンドロイドが人間の生命観に対する新たな視野を提供すること、そして科学と人間性の新たな境界を模索する中でその重要性を理解する。彼はブラック・ジャックが"HeartBeat Mark II"の生命を救うために手術することを支持し、その費用をすぐに提供する。\\
   5. しかし、その手術は革新的な機械科学の知識と技術を必要とし、特に心臓部分の修復は緻密で難解な作業だと分かる。不安に感じつつも、ブラック・ジャックはアンドロイドの問題を解決しようと意志を固める。\\
   6. その間も、ピノコは自らの命を救うために、ドクター・キリコによる拘束から逃れようと全力を尽くしている。\\
   \\
   \{シーン3: 看病・治療\}\\
   1. ブラック・ジャックは川村の信頼と"HeartBeat Mark II"の心臓部分の”生命”を信じて、冷静に複雑な手術に挑む。(ブラック・ジャックの視点)\\
   2. 手術は難関だった。微細な機械構造と専門的な技術が求められる"HeartBeat Mark II"の心臓部分の手術を進めるブラック・ジャック。彼の専門知識の欠如が手術の難易度をさらに上昇させる。\\
   3. しかし、ブラック・ジャックの不屈の努力はやがて結果を伴い、心臓部分の異常は修正され、手術は成功する。\\
   4. ピノコが自力で逃げようとする。その時、ドクター・キリコの部屋からうめき声が聞こえる。\\
   5. 脱出したピノコがドクター・キリコの部屋に来る。ドクター・キリコが胸を押さえて苦しそうにもがいている。そして、気を失う。それを見たピノコがドクター・キリコの異変に驚き、ブラック・ジャックに電話する。\\
   \\
   \{シーン4: 変革\}\\
   1. ブラック・ジャックがドクター・キリコを診察している。ブラック・ジャック「心臓だ！緊急オペを行う。」ピノコが手伝う。\\
   2. 診察台に眠っているドクター・キリコ。ブラック・ジャックはその胸を切開する。その心臓を見て、ブラック・ジャックは驚く。\\
   3. なんと！ドクター・キリコの心臓は、アンドロイドの心臓だった。ブラック・ジャックはそれを手術する。\\
   4. 数日後、ドクター・キリコが目を覚ます。ブラック・ジャックを見て、「何でお前がここに？」ブラック・ジャック「お前に緊急のオペが必要だったんでね。」ドクター・キリコ「何故死なせてくれなかった？」ブラック・ジャック「命を救うのが医者だからな。」ドクター・キリコ「俺はそうは思わない。」ブラック・ジャック「絶対安静だ。しばらく静かにしていろ。安心しろ、オペは確かだから。」と言って去る。ドクター・キリコ「アンドロイドプロジェクトなんて、俺が必ず壊してやる！」と叫ぶ。\\
   \\
   \{シーン5: エピローグ\}\\
   1. 去っていくブラック・ジャックとピノコ。ピノコ「せんせい、本当のことを言わなくていいの？」それに応えず、歩き続けるブラック・ジャック。それについていくピノコ。（終わり）
   
\end{breakitembox}

\vspace{20pt}
\begin{breakitembox}
   [l]{グループCが提出したプロット}
   \noindent\{シーン1: 出会い\}\\
   1. 自宅にて気鬱とした日々を送っていたトシロウの元に、一匹の野良猫が現れた。\\
   2. その小さくて愛らしい姿に触れ、トシロウは野良猫にミーちゃんと名前をつけた。\\
   3. ミーちゃんとの出会いを通して気分が晴れやかになったトシロウは、どこに行くにもミーちゃんを連れ歩いた。その行動の一部にはブラック・ジャックの診療所への訪問も含まれていた。\\
   4. ブラック・ジャックの診療所で、ミーちゃんはピノコに興味を持ち、彼女の手に自分の小さな爪をタッチさせた。\\
   5. ミーちゃんとトシロウの関係は日に日に深まり、トシロウはミーちゃんに、鈴の付いた首輪をプレゼントした。\\
   \\
   \{シーン2: サポート\}\\
   1. トシロウを定期的に見舞う孫のケイコが訪れ、高齢で療養中のトシロウがミーちゃんを飼うことに反対し、口論になる。\\
   2. ミーちゃんはケイコの言葉に打ちのめされ、トシロウの前から姿を消したが、彼の病状を遠くから見守り続けていた。\\
   3. ミーちゃんがいなくなり、町中を探し歩くトシロウの後姿を、偶然ピノコが目撃した。\\
   4. 一方、ケイコはミーちゃんに新たな里親を探すことを決め、町中の動物愛護団体を訪問し始めた。\\
   5. トシロウは、空腹のミーちゃんを気遣い、心の中でミーちゃんに話しかけながら毎晩門の外にエサを置く。翌朝、きれいになった皿を見て、トシロウは希望を持ち続けた。\\
   \\
   \{シーン3: 危害\}\\
   1. 数日後、トシロウは心労が重なり、自宅で倒れてしまった。それを見たミーちゃんは一旦部屋の中に入り、すぐにブラック・ジャックの診療所へ向かった。\\
   2. ピノコは、ミーちゃんが持ってきたトシロウの診療カードを見つけ、彼の危険な状態を察知し、大声でブラック・ジャックを呼んだ。\\
   3. ブラック・ジャックは、ミーちゃんが暑い中、短時間で多くの移動を経験した結果、体調を崩していることを見抜いた。\\
   4. 同時に、ブラック・ジャックは自分の患者のケアが充分でなかったことを後悔し、直ちに往診の準備を始めた。\\
   5. 一方、ミーちゃんの体調はさらに悪化し、振るえが止まらず、とうとう意識を失ってしまった。\\
   \\
   \{シーン4: 依頼\}\\
   1. ブラック・ジャックはピノコに状況を説明し、対応を任せ、自らはトシロウのもとへと急行した。\\
   2. ブラック・ジャックが訪れると、トシロウは意識混濁の中、ミーちゃんの名を呼び続けていた。\\
   3. ブラック・ジャックは、トシロウの容態の悪化は持病だけでなく、ミーちゃんの不在によるものであると理解した。\\
   4. ブラック・ジャックは、トシロウの抜本的な治療を行うために、彼を自身の診療所へと搬送することに決めた。\\
   5. ブラック・ジャックは、トシロウを診療所に連れて行くことをメモに残した。\\
   \\
   \{シーン5: 看病・治療\}\\
   1. 診療所に到着したブラック・ジャックは、全力でトシロウの治療に取り組んだ。\\
   2. 一方で、ピノコは別室ですっかり弱ったミーちゃんの看病を優しく行い、回復に寄り添った。\\
   3. お互いにそれぞれ尽力したおかげで、トシロウとミーちゃんの体調は徐々に回復の兆しを見せ始めた。\\
   4. そしてトシロウに、ミーちゃんが隣の部屋にいると告げた途端、彼の表情は一転、生命力が湧いてきた。そのとき、トシロウはミーちゃんがつけているのと同じ鈴が離れた場所から聞こえるのを捉えた。\\
   5. ミーちゃんも自身の健康を取り戻しつつあり、再びトシロウと対面することを願い出始めた。\\
   \\
   \{シーン6: 反省と理解\}\\
   1. ミーちゃんを引き取りにきた保護団体から、トシロウの不在を知らされたケイコは、トシロウの家を訪れ、ブラック・ジャックから残されていたメモを発見した。\\
   2. 彼女はそのメモの指示に従って診療所を訪れ、ミーちゃんとトシロウが同じ場所で治療を受けている事実に激怒した。\\
   3. ブラック・ジャックからトシロウとミーちゃんの絆の深さを諭されたケイコだが、心には響かなかった。\\
   4. しかし、診療所を後にしたケイコは道すがら、トシロウが小さな自分の髪につけてくれた鈴を思い出した。\\
   5. ケイコはその記憶とミーちゃんに付けられた鈴とを重ね合わせ、今のトシロウにとってミーちゃんがかけがえのない存在だと理解し、同時に今でも自分がトシロウから愛されていると感じた。\\
   \\
   \{シーン7: 感謝\}\\
   1. ついに全快したトシロウは、退院する日にミーちゃんと念願の再会を果たし、ブラック・ジャックは彼の回復を喜んだ。\\
   2. ピノコは、ミーちゃんの優しさと強さを尊敬し、自分自身の頑張りも誇らしく思い、ブラック・ジャックを見上げた。\\
   3. そのときドアが開き、笑顔のケイコが入ってきた。ケイコはミーちゃんとトシロウが幸福そうに過ごす様子を見て涙があふれ、ブラック・ジャックに深々と頭を下げた。\\
   4. ブラック・ジャックは今回の報酬について「猫のミーちゃんとトシロウの間に芽生えた絆は、有り金全てで購入したいほどの価値がある」と言った。\\
   5. ケイコはそっとミーちゃんを抱き上げ、三人は笑顔で手を振りながら診療所を出て行き、そこには家族の温かさが溢れ出ていた。\\
   6. その後ろ姿を見送るブラック・ジャックとピノコの元に、ケイコがトシロウの家に引っ越すことを決めたという声が遠くから聞こえてきた。その言葉を聞いた二人は互いに目を合わせ、満面の笑みを浮かべる。ブラック・ジャックはピノコの頭をなで、「良い選択だ」と目を細めて呟いたのだった。
   
\end{breakitembox}

\vspace{20pt}
\begin{breakitembox}
   [l]{グループDが提出したプロット}
   \noindent シーン1: 再会\\
   1.（主人公視点）徹底した環境活動家、如月めぐみは、冷酷な製薬王・白根至の鉄壁の研究所で、彼女と因縁の深い天才医師ブラック・ジャックと運命的に再会する。\\
   2.彼女は、ブラック・ジャックとその無邪気な助手ピノコに向き合い、地球に迫りくる破滅の危機を千切れる声で叫ぶ。\\
   3.その危機は、無慈悲な白根至が自然から掠め取った成分で開発した新薬による母なる大地の創造力の破壊であり、めぐみはそれに孤独で果敢に立ち向かう。\\
   4.ブラック・ジャックは、死闘を繰り広げる彼女の献身に心を動かされ、地球に命のために桁外れの力を尽くす彼女を陰から見守ることを決める。\\
   5.如月の行動が地球の運命と密接に結びつき、それが一人の生物として地球の存続を壮大なスケールで左右する。\\
   \\
   シーン2: 警告\\
   1.（協力者視点）ブラック・ジャックとピノコは、如月から聞かされた自然界の破壊の現状に心底震える。\\
   2.彼らは、自分達の雇い主である白根至の新薬開発が地球への大きな脅威になりつつあることを肌で感じる。\\
   3.ブラック・ジャックとピノコは、如月と共にこの壮絶な戦いに身を投じる決意を固める。\\
   4.ブラック・ジャックは、抗うべき自然破壊の過程で、自身の医術も新たな進化の可能性を探ることを心に決意する。\\
   5.如月と共に、ブラック・ジャックは自然破壊の防止の叫びを巻き起こし、未知なる挑戦への道を進む。\\
   \\
   シーン3: 情報開示\\
   1.（協力者視点）ブラック・ジャックとピノコは、自然の成分の詳細な調査を徹底的に行い、白根至の新薬の危険性を発見する。\\
   2.ブラック・ジャックは、白根至の秘密裏に進行していた新薬開発の裏側に仕組まれた陰謀を突き止める。\\
   3.ブラック・ジャックとピノコは、その情報を手に入れ、白根至の意図を真摯に追求し、地球の防衛につながる道を見つける。\\
   4.ブラック・ジャックは自然の力を活かした治療法を開発し、如月の協力を得て新薬の裏側に隠された真実を明らかにする。\\
   5.地球の危機が一層深刻化する中、彼らは決して諦めず、地球の防衛を新たに誓う。\\
   \\
   シーン4: 発見\\
   1.（協力者視点）ブラック・ジャックは、如月から提供される情報を基に、白根至の新薬の裏側に隠された真実を解き明かす。\\
   2.その真実は、彼らに衝撃を与え、新たな闘争の火種を作り出す。\\
   3.ブラック・ジャックは、医学と生態学を結びつけ、如月の危機に立ち向かうための革新的な治療法を提案する。\\
   4.彼とピノコは、自然と科学の力を結合した新たな治療法、 ”エコメディカル”を開発し、如月の身に忍び寄る危険に対抗する。\\
   5.地球の窮地を前に、彼らは白根至との闘いでも一定の勝利を得る。\\
   \\
   シーン5: 危機\\
   1.（協力者視点）如月とブラック・ジャックは、地球の存続を訴え、絶対的な困難にも関わらず行動を共にする。\\
   2.しかし、彼らが地球を救おうと絶え間なく活動する一方で、白根至は一層強固な抵抗を展開する。\\
   3.未来を見据えた如月とブラック・ジャックは、決して諦めず新たな挑戦へと突き進む。\\
   4.物語は、彼らが全てを乗り越え地球を救うことができるのか、そして白根至のあくどい陰謀から逃れられるのかという緊張感に満ちたクライマックスに向けて突進する。\\
   5.物語は、未来への再生の希望とともに終わり、ブラック・ジャック、如月、そして地球の未来が次章で明らかになる。
\end{breakitembox}

\vspace{20pt}
\begin{breakitembox}
   [l]{グループEが提出したプロット}
   \noindent \# シーン1: 挑発\\
   1. 場刷律という人物が、美澄紺という名の女優志望の若者を次回のビデオ撮影にスカウトするという情報がピノコから伝えられてきた。\\
   2. SNS上で挑戦的な言動を振りまく彼が自分のYouTubeチャンネル視聴者を増やす為に、紺を利用しようとしている事が見て取れた。\\
   3. それでも美澄紺は、やや彼の挑発には乗る形を見せ、リスキーな撮影を受け入れてしまっている。\\
   4. 紺自身、律の目的を理解していないことが明らかで、これが自身の女優キャリアを高める手段だと考えているようだ。\\
   5. 彼女は自身がメディア上で注目を集めることが重要だという考えを持っているようだ。\\
   \\
   \#\# シーン2: 対立\\
   1. 撮影が続く中で、限界を超えた要求がされ、紺は抵抗をし始めた。\\
   2. 彼女は自己防衛と撮影の安全性の意識から場刷律に対して撮影のリスクについて意見し始めた。\\
   3. しかし、場刷律は紺の懸念を無視し、視聴者数を増やすためにはリスクを取るべきだと主張する。\\
   4. こうした齟齬から二人の間には深刻な対立が生まれる。\\
   5. 場刷律は彼女の言葉に反発し、自分の理想を守るために対立していくことを決意する。\\
   \\
   \#\# シーン3: 危害\\
   1. ブラックジャックは呼ばれて駆けつけ見たのは、撮影に夢中だった律の無茶な要求により紺が事故に遭い、倒れている姿だった。\\
   2. 生配信され、視聴者たちが動揺する中、紺はうめきながら倒れていった。\\
   3. その一部始終は律のカメラに捉えられており、SNSでは大騒ぎになっていた。\\
   4. 律は紺の事故により多くの注目を集めてしまったが、その一方で視聴者数も増えてしまっている事実を目の当たりにしながら。\\
   5. 律はブラックジャックに治療を頼むべく、紺を診療所へと運んできた。\\
   \\
   \#\# シーン4: 叱責\\
   1. ブラックジャックとピノコが紺の深刻な状況を見て息を呑んだ。\\
   2. ブラックジャックは紺の治療に必要な準備を整え、すぐに手術を開始した。\\
   3. 彼は手術中にも場刷律を叱責し、彼の行いが人の命を軽んじる行為だと説明した。\\
   4. ピノコも泣きながら紺の手を握り、彼女が回復することを祈っていた。\\
   5. 場刷律はブラックジャックからの叱責を受け止め、人間の命の重さを再認識する。\\
   \\
   \#\# シーン5: 発見\\
   1. 手術の結果、紺は無事に意識を取り戻した。\\
   2. ブラックジャックとピノコは安心し、紺にはゆっくりと休んで欲しいと伝え、紺の看病を続けた。\\
   3. 律は紺の本当の目的を理解するが同時に不満も抱いていた。\\
   4. 紺は彼の行為が自分の女優としての努力を無駄にしてしまうかもしれないと気づいていた。\\
   5. さらに、律のせいで自身の信用が傷つき、女優として活動していくことが難しくなることに気づく。\\
   \\
   \#\# シーン6: 真相発覚\\
   1. 紺は自身の回復を果たし、再び人々の前に立つ決心を固めていた。\\
   2. そして彼女は、律が自己中心的な言動の事実を視聴者たちに開示すると決意していた。\\
   3. その為に彼の自己中心的な行動をSNSで暴露し、視聴者たちに事実をまざまざと見せつけた。\\
   4. 彼女の行動に対し多くの視聴者が支持を示し、一方で律は多くの視聴者やフォロワーを失っていった。\\
   5. 律は自分の過ちについてSNS上で謝罪する一方で、紺は再起し女優としての人生を歩んで行く決意を新たにしていた。
\end{breakitembox}


\expandafter\ifx\csname ifdraft\endcsname\relax
   \fancyhead[LO]{}
   \printbibliography[title=参考文献]
   \end{document}
   \fancyhead[LO]{\rightmark}
   \fi

